% SEAMDB -- design and rationale document.
% Compile on Overleaf (pdfLaTeX). Only ubiquitous packages are used:
% geometry, xcolor, hyperref, booktabs, listings. No shell-escape, no minted.
% Source is ASCII-only on purpose (matches the repository lint rule).

\documentclass[11pt]{article}

\usepackage[margin=1in]{geometry}
\usepackage[T1]{fontenc}
\usepackage{lmodern}
\usepackage{microtype}
\usepackage{xcolor}
\usepackage{booktabs}
\usepackage{listings}
\usepackage{hyperref}

\hypersetup{
  colorlinks=true,
  linkcolor=blue!50!black,
  urlcolor=blue!50!black,
  pdftitle={SEAMDB: a proof-carrying evidence plane over a from-scratch relational engine},
  pdfauthor={Mitrajit Ghorui},
}

\definecolor{codebg}{RGB}{246,246,244}
\definecolor{codekw}{RGB}{16,92,168}
\definecolor{codecm}{RGB}{96,128,96}

\lstdefinelanguage{SEAMQL}{
  morekeywords={spawn,make,table,drop,in,add,index,on,let,allow,block,deny,
    tweak,update,yeet,remove,pull,read,from,left,inner,link,join,when,where,
    bunch,group,by,rank,sort,cap,take,start,begin,tx,lockin,commit,dip,cancel,
    spill,explain,main,key,needed,links,to,whole,number,text,decimal,bool,
    stamp,day,as,and,or,not,is,null,check,ship,true,false},
  sensitive=true,
  morecomment=[l]{--},
  morestring=[b]',
}
\lstset{
  language=SEAMQL,
  basicstyle=\ttfamily\small,
  keywordstyle=\color{codekw}\bfseries,
  commentstyle=\color{codecm}\itshape,
  stringstyle=\ttfamily,
  backgroundcolor=\color{codebg},
  frame=single,
  rulecolor=gray!40,
  columns=flexible,
  keepspaces=true,
  showstringspaces=false,
  breaklines=true,
  aboveskip=0.8em,
  belowskip=0.8em,
}
\newcommand{\seamql}{\lstinline[language=SEAMQL]}
\newcommand{\code}[1]{\texttt{#1}}

\title{SEAMDB: A Proof-Carrying Evidence Plane\\Over a From-Scratch Relational Engine}
\author{Mitrajit Ghorui \\ Genauix}
\date{October 2026}

\begin{document}
\maketitle

\begin{abstract}
Regulated data handoffs (model exports, claims extracts, cross-org feature
feeds) fail today for an evidentiary reason rather than a storage reason: the
receiver cannot prove \emph{what} left, \emph{under which policy}, from
\emph{which source snapshot}, and whether the sender later rewrote history.
SEAMDB answers that with an \emph{Evidence Capsule}: a canonicalized,
content-addressed bundle -- source snapshot roots, lineage closure, policy
decision digest, execution appraisal, result root, and a transparency-log
receipt -- that an auditor verifies \emph{offline}, where any mismatch exits
nonzero. Underneath, SEAMDB is also a complete relational engine written from
scratch in TypeScript (typed pages, CRC-checked WAL, snapshot-isolation MVCC,
B$+$tree, and a query language called SEAMQL), because the evidence plane must
own the write path to fingerprint statements, plans, and write sets. This
document describes the architecture end to end, justifies each design choice
against the obvious alternatives (Redis, PostgreSQL, SQLite, MongoDB, DuckDB,
audit tables, external transparency logs, ad-hoc export scripts), states the
advantages honestly, and lists the limitations without decoration. The whole
system is verifiable: 204 engine tests, a ten-attack adversarial gallery, and
known-answer vectors all run offline. The same engine serves clients over a
framed wire protocol: a five-page web site (home, SQL IDE, SECP lab, specs,
docs) that runs it in the browser, and a native PyQt6 desktop IDE (SEAMDB
Studio) with a schema explorer, syntax-highlighted editors and a results
grid, driven by a Python client. The wire layer is one protocol (frame types
1--8: hello, auth, query, result, error, cancel, status, catalog) with a
length-prefix transport; both languages bind to identical semantics. The
Python smoke run and the catalog introspection are covered by the test
suite, and the desktop IDE ships a Qt syntax highlighter kept in sync with
the canonical SEAMQL keyword surface.
\end{abstract}

\tableofcontents

% =====================================================================
\section{The problem: exports without custody}
\label{sec:problem}

An export pipeline in a regulated setting is a custody chain: source systems
of record, a transformation, a policy decision, a destination. Today that
chain is documented, not \emph{evidenced}. When an auditor asks ``prove that
this file came from the March claims snapshot, under policy bundle 42, and was
not modified afterwards,'' the answer is a pile of logs that must be trusted
wherever they were written. Three concrete failure modes motivate SEAMDB:

\begin{enumerate}
  \item \textbf{Silent rewrite.} The sender can regenerate the export from
    different inputs and the receiver has no commitment that pins the
    original. Hashing the final file only proves the file was not changed in
    transit -- it says nothing about \emph{which inputs} produced it.
  \item \textbf{Policy theater.} Approvals live in tickets and wikis, not in
    the artifact. Nothing stops the next export from ignoring the decision,
    and a stale approval looks identical to a fresh one.
  \item \textbf{Evidence that only the author can check.} Most ``audit''
    features require calling back into the producer's systems, so the party
    under investigation is also the party running the verification.
\end{enumerate}

SEAMDB's thesis: \emph{every high-risk export must carry its own
verification, checkable by anyone, without calling home.}

% =====================================================================
\section{What SEAMDB is}
\label{sec:what}

SEAMDB has two coupled halves:

\begin{itemize}
  \item An \textbf{evidence plane}: canonicalization, keyed Merkle
    commitments, a transparency log with signed checkpoints, witness
    acknowledgements, a policy gate, an execution appraisal, one-use release
    grants, and the Evidence Capsule itself -- a single wire object whose
    digest covers every sealed field.
  \item A \textbf{relational engine written from scratch} (SEAMDB Engine):
    16\,KiB pages with CRC32 on every read, a write-ahead log with
    ARIES-style analysis/redo/undo recovery, snapshot-isolation MVCC, a
    B$+$tree with unaligned slot pages, a buffer cache, page-level
    AES-256-GCM encryption, and SEAMQL -- a typed query language with
    lexer, parser, resolver, and planner implemented in this repository.
\end{itemize}

The engine is not a vanity project: the evidence plane must be able to say
\emph{which statement, which plan, which write set} produced a result, and
that is only honest if the database itself emits those fingerprints. A
Postgres adapter could be added later as a capture connector; the built-in
engine is the reference implementation and the browser-runnable demo.

SEAMDB is an \emph{overlay}, not a replacement warehouse: source systems stay
systems of record; capture is read-only; raw payloads stay in
customer-controlled storage. What crosses the boundary is the capsule.

% =====================================================================
\section{Architecture}
\label{sec:architecture}

\begin{verbatim}
  customer data plane             SEAMDB control plane
  (Postgres, warehouse,    -->   intent / snapshot / seal /
   feature store: SoR)            grant / verify (in process)
          |                                |
          | read-only capture              | SECP/0.1 state machine
          v                                v
  +----------------------------------------------------------+
  | evidence plane: JCS canon -> keyed Merkle -> capsule     |
  | digest -> transparency log -> witnesses -> verifier      |
  +----------------------------------------------------------+
          |                                |
          v                                v
  trust plane (split duty)         consumer plane (auditor)
  log checkpoints, witnesses,      capsule + npm run verify,
  offline verifier CLI             exit 0 only on full PASS
\end{verbatim}

\subsection{Data plane}
Sources register a schema and a connector version; capture snapshots rows
under a keyed Merkle tree (HMAC-SHA-384 leaves, domain-tagged
\code{SEAMDB/TLOG-LEAF/V1}-style per-purpose tags). Identical snapshot tokens
produce identical roots, so a re-capture is provably the same data.

\subsection{Control plane}
An in-process API with SECP/0.1 semantics: \code{intent}, \code{snapshot},
\code{seal}, \code{grant}, \code{verify}. Illegal state transitions throw;
\code{DENIED} is absorbing except from \code{RELEASED}. The operator may
authorize a release but cannot mutate a sealed snapshot without changing its
digest -- the enforcement is algebraic, not procedural.

\subsection{Evidence plane}
RFC 8785 JSON Canonicalization Scheme (JCS) produces the byte string that
gets hashed; pretty-JSON hashing is treated as a defect. Every commitment is
SHA-384 with an explicit domain tag; leaves are keyed (HMAC) so an adversary
cannot grind alternative leaf preimages without the commitment secret. The
capsule digest is
\[
  \mathrm{SHA384}\big(\texttt{SEAMDB/CAPSULE/V1}\,\|\,\mathrm{JCS}(\mathrm{body})\big),
\]
which covers policy purpose and destination, execution profile, result root,
row count, key epoch, and the transparency pointers. Mutating any sealed
field changes the digest -- this is enforced by tests that mutate each field
in turn and assert the digest moves.

\subsection{Trust plane}
Checkpoints over the log are signed with the runtime's classical signature
algorithm -- Ed25519 where WebCrypto provides it, ECDSA P-384 (FIPS 186-5)
otherwise -- and \emph{every checkpoint names which algorithm produced it}.
Witness acknowledgements (quorum 2 in the demo) bind witnesses to
checkpoint roots; \code{detectSplit} flags equivocation: same tree size,
different root. The verifier CLI is a sibling script of the attack gallery,
so the party checking evidence is not the party that produced it.

\subsection{Consumer plane}
The auditor receives the capsule plus proofs and runs \code{npm run verify}
offline. Exit code 0 only on a full PASS across all named checks;
\code{INDETERMINATE} is explicitly not a release signal. There is no network
access in the verifier path at all.

% =====================================================================
\section{Storage engine}
\label{sec:storage}

\subsection{Pages and encoding}
Pages are 16{,}384 bytes (a power of two that fits NVMe and object-store
multipart boundaries), each prefixed by a header and protected by CRC32
verified on \emph{every} read -- a corrupted page fails closed with a typed
code (\code{E\_CORRUPT}) rather than serving damaged rows. Slot pages use
unaligned variable-length records, so text and decimals do not pay for
maximum-width padding. Values are typed: \emph{whole number} (int64),
\emph{decimal(p,s)} (exact, no binary float), \emph{text}, \emph{bool},
\emph{stamp}, \emph{day}. There is no implicit numeric conversion; a type
mismatch raises \code{E\_TYPE\_MISMATCH} at resolve time.

\subsection{Write-ahead log and recovery}
Every commit is durable in the WAL before acknowledgment. The WAL itself is
CRC-framed; recovery performs ARIES-style analysis, redo of committed work,
and undo of loser transactions, then writes a checkpoint. The crash matrix
tests tear the log at adversarial byte offsets (including split frames) and
require byte-exact logical recovery or a clean fail-closed refusal --
never a partially applied transaction.

\subsection{MVCC}
Readers operate on snapshots (transaction id plus an active-transaction set);
writers create new versions instead of overwriting. Consequences that the
IDE demonstrates live: an explicit transaction sees its own writes, other
sessions never see uncommitted rows, and \seamql{dip} (abort) makes every
write vanish -- the demo script inserts inside a transaction, reads it back,
aborts, and shows the row gone. Locking is per main-key with a deadlock-free
acquisition order; a second \seamql{dip} without a transaction raises
\code{E\_TX\_STATE}.

\subsection{B$+$tree and planner}
Secondary indexes are B$+$trees over typed keys. The planner is honest about
what it does: \seamql{spill} (explain) prints the physical plan, and the
SQL IDE page snippet shows the plan flip from
\code{HeapScan(people) -> Filter -> Project} to
\code{IndexPoint(people.people\_by\_id) -> Filter -> Project} after
\seamql{add index} -- the same query, the same rows, a different access path
the user can see.

\subsection{Encryption at rest}
Page-level AES-256-GCM: data keys are derived per object from a tenant
secret via HKDF-SHA-384 with a key-epoch salt, and every page's additional
authenticated data binds tenant, object, type, schema version, and epoch.
A page moved between objects or epochs fails authentication instead of
decrypting. KMS/HSM custody is \emph{simulated} (envelope keys over real
AES-256-GCM) in this build and labeled as such everywhere it appears.

% =====================================================================
\section{SEAMQL}
\label{sec:seamql}

SEAMQL is a small typed language designed for auditability: every construct
maps to one engine operation, and every operation can fail only in ways the
error catalog enumerates. A representative script:

\begin{lstlisting}
-- DDL: catalog write, WAL-cat
spawn table folks: id: whole number, main key,
                    name: text, needed
spawn table runs:  rid: whole number, main key,
                    folks_id: whole number, needed,
                    links to folks.id,
                    amount: decimal(12, 2), needed

drop in folks { id: 1, name: 'Ada' }
drop in folks { id: 2, name: 'Grace' }
drop in folks { id: 3, name: 'Linus' }
drop in runs  { rid: 1, folks_id: 1, amount: 10.00 }

-- LEFT JOIN: unmatched customers null-extend
pull folks.name as customers
  from folks
  left link runs on folks.id = runs.folks_id
  when runs.rid is null
\end{lstlisting}

This is LeetCode 183 (``Customers Who Never Order''), executed verbatim by
\code{tests/demo-proof.test.ts} and by the browser IDE on the SQL IDE page
(/ide.html), returning \code{Grace, Linus}. The language surface: \seamql{spawn} (DDL),
\seamql{drop in} (insert), \seamql{tweak} (update), \seamql{yeet} (delete),
\seamql{pull} with \seamql{left link}/\seamql{join}, \seamql{when} filters,
\seamql{bunch by} (group), \seamql{rank by} (order), \seamql{cap} (limit),
\seamql{add index}, \seamql{start tx} / \seamql{lockin} / \seamql{dip},
\seamql{spill} (explain), and \seamql{let}/\seamql{block} row policies.
Constraints are enforced by the engine, not the application: re-inserting a
main key raises \code{E\_UNIQUE}; a dangling foreign key raises \code{E\_FK}.

\subsection{Errors are typed}
Every failure carries a stable code and a source position: \code{E\_UNIQUE},
\code{E\_FK}, \code{E\_UNKNOWN\_NAME}, \code{E\_TYPE\_MISMATCH},
\code{E\_TX\_STATE}, \code{E\_LIMIT}, \code{E\_CORRUPT},
\code{E\_WAL\_CORRUPT}, \code{E\_DENIED}, \code{E\_CRYPTO},
\code{E\_FEATURE\_NOT\_SUPPORTED}. The IDE renders them per statement and
keeps executing subsequent statements, so a demo script can show three
different typed failures in one run -- which is exactly what the ``Fail
closed'' snippet on the SQL IDE page does.

% =====================================================================
\section{Evidence plane in detail}
\label{sec:evidence}

\subsection{Canonicalization}
JCS fixes one byte representation per JSON value (sorted keys, no
whitespace, shortest-number form). All digests in the system are taken over
canonical bytes or over explicitly domain-tagged byte strings -- never over
``whatever \code{JSON.stringify} emitted this run.'' The canonicalizer is
covered by known-answer vectors.

\subsection{Merkle profile}
Leaves are \code{H(0x00 || tag || keyed(row))}, nodes are
\code{H(0x01 || tag || left || right)}, empty trees hash a dedicated empty
tag, and odd nodes promote (RFC 9162-shaped). Inclusion and consistency
proofs are verified by the offline verifier; the known-answer suite pins
tree shapes (including 0, 1, and odd leaf counts).

\subsection{Transparency log}
Every sealed capsule appends a leaf; each append emits a signed checkpoint
plus an inclusion proof and, when applicable, a consistency proof.
Equivocation detection compares local and remote checkpoints at equal tree
sizes. The demo witnesses are simulated (HMAC-based acknowledgements), and
the documentation says so -- the \emph{logic} of quorum checking and split
detection is real; the federation is not.

\subsection{The nine-gate filtration (SECP/0.1)}
\code{INTENT\_CREATED} $\rightarrow$ \code{SOURCE\_COMMITTED} $\rightarrow$
\code{LINEAGE\_CLOSED} $\rightarrow$
\code{POLICY\_APPRAISED} $\rightarrow$ \code{EXECUTION\_APPRAISED}
$\rightarrow$ \code{RESULT\_COMMITTED} $\rightarrow$
\code{CAPSULE\_SEALED} $\rightarrow$ \code{RECEIPT\_VERIFIED}
$\rightarrow$ \code{RELEASED}, with \code{DENIED} absorbing. Each gate
corresponds to a named check in the verifier; an illegal transition throws.

\subsection{Policy and grants}
The policy decision binds purpose, destination, data classes, query profile,
row/byte caps, an expiry, a one-use nonce, and the policy bundle digest. A
stale bundle digest (\code{sha384:dead} in the demo) or an open critical
discrepancy denies the grant. Grants are one-use and bound to the capsule
digest: replaying a consumed grant fails closed.

\subsection{Coherence and algebraic checks}
Discrepancy witnesses compare overlapping views (same \code{member\_ref}
key) and \emph{retain} the residual rather than silently resolving it; an
algebraic conservation check asserts that input sums equal output sums over
the overlap. An MDL-style compression residual scores anomaly-ness and is
explicitly labeled ``not Kolmogorov complexity $K(x)$.''

\subsection{The verifier}
\code{npm run verify} runs every named check over an exported bundle with no
network access: capsule digest recomputation, inclusion and consistency
proofs, witness quorum, policy binding, appraisal, conservation, and grant
state. Verdicts are \code{PASS}/\code{FAIL}/\code{INDETERMINATE}; anything
but full PASS exits nonzero. The adversarial gallery seeds at least ten
attacks (leaf tamper, purpose swap, policy stale-roll, grant replay,
equivocation, downgrade, \ldots) and the suite asserts each is blocked.

% =====================================================================
\section{Cryptography profile}
\label{sec:crypto}

All primitives are imported, never invented; the repo's injection-surface
constraints (Section~\ref{sec:constraints}) forbid \code{eval}, dynamic
\code{require}, and WebAssembly outright.

\begin{center}
\begin{tabular}{@{}llp{7.6cm}@{}}
\toprule
Purpose & Primitive & Notes \\
\midrule
Commitments & SHA-384 (WebCrypto) & domain-tagged per use \\
Keyed leaves & HMAC-SHA-384 & commitment secret never leaves process \\
Encryption & AES-256-GCM & 12-byte nonce, AAD binds object + epoch \\
Key derivation & HKDF-SHA-384 & per-object DEK, epoch salt \\
Transport & ML-KEM-768 + X25519 & hybrid, FIPS 203 via @noble/post-quantum \\
Transport fallback & ML-KEM-768 + ECDH P-256 & where WebCrypto lacks X25519 \\
Checkpoints & Ed25519 or ECDSA P-384 & per-checkpoint algorithm label \\
\bottomrule
\end{tabular}
\end{center}

\subsection{Hybrid transport, honestly labeled}
The transport handshake mirrors the TLS~1.3 \code{X25519MLKEM768} hybrid
group: both an ephemeral classical agreement and ML-KEM-768 encapsulation run
per session, and the KEM shared secret is a \emph{mandatory} HKDF input --
a classical-only peer derives a different traffic secret, so the handshake
cannot be silently downgraded. Both encapsulation keys (by SHA-384 digest)
and the agreement group are folded into the \code{info} string. Every
negotiation record names the groups that actually ran; on runtimes whose
WebCrypto lacks X25519 (some browser builds), the engine falls back to
ECDH P-256 and says so in the profile label and note. The fallback choice
comes from probing the local runtime only -- never from wire input -- so an
attacker cannot force it.

\subsection{Signature agility with labels}
Likewise, checkpoint signatures use Ed25519 where available and ECDSA
P-384 otherwise, and every checkpoint carries an \code{alg} field. The
product rule behind both mechanisms: \emph{the record says what happened};
the UI never prints a stronger claim than the runtime executed. This is why
the site can promise ``hybrid PQ transport'' and ``classical
signatures'' on a browser build that lacks X25519/Ed25519 without lying --
it runs P-256 and labels it.

\subsection{What the crypto does not claim}
PQ is transport only: there are no PQ signatures, and ML-KEM never
``attests'' a capsule. The system is not a custom cipher, not a TLS stack,
and not a zero-knowledge proof system. Timestamps in the demo are host-local
(a real deployment needs a trusted timestamping source).

% =====================================================================
\section{Why not the other database?}
\label{sec:alternatives}

The most common reactions to SEAMDB, and the honest answers.

\subsection{``Just use Redis.''}
Redis is an excellent in-memory data structure server and a superb cache and
broker. It was never asked to produce export evidence: keys hold opaque
values the server does not interpret; persistence (AOF/RDB) trades durability
windows for speed by design; there is no snapshot isolation between
commands; there is no query language with joins, types, and constraints; and
nothing in the trust story lets an auditor verify an export offline. The
question is not ``which stores data faster'' -- Redis wins that -- but
``which store's writes carry proof.'' Cache vs.\ custody. Keep Redis as a
cache; it does not become evidence infrastructure by adding a hashing
wrapper.

\subsection{``Just use PostgreSQL.''}
Postgres is the correct default for almost everything, and SEAMDB assumes it
as a system of record. But the evidence question lives \emph{above} any
single database: Postgres gives you MVCC, WAL, and logical replication --
all excellent -- and no capsule. A logical-decoding-based exporter in front
of Postgres can capture change streams, but then the evidence plane must
still be built: canonicalization, keyed commitments, policy binding,
transparency log, offline verifier. That is precisely the layer SEAMDB
\emph{is}. Using Postgres for storage would trade a from-scratch engine
(Section~\ref{sec:limitations}: single writer, no parallel query) for a
dependency whose internals the capsule cannot fingerprint end to end. Both
positions are defensible; SEAMDB chose the one where every fingerprint --
statement, plan, write set -- originates in code we can show.

\subsection{``Just use SQLite.''}
SQLite is a marvel of engineering: embeddable, battle-tested, and its WAL
mode is excellent. It also has no multi-writer story for our target
deployment, no network protocol by design (which is fine -- SEAMDB's
protocol is also in-process), and, decisively, no evidence plane: its
\code{dbstat}/dump artifacts prove shape, not custody. SEAMDB's engine is in
the same embeddable spirit but adds the parts SQLite deliberately excludes:
typed constraint enforcement with typed errors, keyed commitments per
snapshot, statement/plan digests in the log, and a capsule format.

\subsection{``Just use MongoDB (or any document store).'''}
Document stores optimize schema flexibility; SEAMDB optimizes schema
\emph{evidence}. A capsule must bind a declared schema digest so that a
drifted schema is detectable; flexible-schema systems make that claim
vacuous. Type coercion rules in document stores also undermine conservation
checks (the algebraic sum-equality over overlaps needs exact numerics --
hence \emph{decimal(p,s)}, not IEEE floats or ``whichever JSON number came
through'').

\subsection{``Just use DuckDB (or any OLAP engine).'''}
Analytical engines answer queries over columnar data quickly; they do not
publish custody. DuckDB inside the transformation step would be entirely
compatible with SEAMDB -- indeed the natural place for heavy transforms --
because SEAMDB's claim is about the \emph{boundary} (what left, under what
policy, provable later), not about compute speed. If your only problem is
fast analytics, do not add SEAMDB.

\subsection{``Just add triggers and an audit table.''}
Append-only audit tables are the most common real-world answer, and they
fail the rewrite test: an operator with write access can rewrite the audit
table too, and both tables will agree afterwards. Evidence requires a
\emph{digest chain plus an append-only log whose checkpoints are signed and
witnessed}, and verification must be runnable by the counterparty. SEAMDB is
that, packaged: transparency log, witness acks, split detection, and an
offline verifier whose exit code is the verdict.

\subsection{``Just run a transparency log (or Sigstore) on the side.''}
External transparency logs solve append-ordering and equivocation detection
well, and SEAMDB's log is deliberately RFC-9162-shaped so a production
deployment can \emph{replace} the local log with a public one. But a log
proves \emph{that a statement was logged}, not that the statement is true:
the log does not know that the capsule's result root came from the claimed
source snapshot under the claimed policy. SEAMDB's contribution is the
binding -- source snapshot, lineage, policy decision, appraisal, result --
that makes the logged statement mean something. Logs and bindings are
complementary; SEAMDB ships both and keeps the seam explicit.

\subsection{``Just hash the CSV on export.''}
This is the ``trust me'' answer in miniature: a file hash proves integrity
since hashing, provenance not at all. It does not bind the inputs, the
policy, the operator, or the time; it cannot detect that the same file was
regenerated from different rows; and it cannot be checked without the
sender's cooperation for context. The capsule is what ``hash the export''
wants to be when it grows up.

\subsection{Summary table}
\begin{center}
\small
\begin{tabular}{@{}p{2.9cm}p{3.6cm}p{3.4cm}p{3.6cm}@{}}
\toprule
& Wins at & Does not answer & SEAMDB's stance \\
\midrule
Redis & microsecond cache/broker & durability windows; no typed queries; no
  proofs & keep it as cache; not custody \\
Postgres & general-purpose SoR, replication & no capsule, no offline
  verification & assume as source; evidence layer sits above \\
SQLite & embeddable single-file DB & no multi-writer; no evidence plane &
  same embeddable spirit, plus fingerprints \\
MongoDB & schema flexibility & schema-drift detection; exact numerics &
  schema digest is part of the capsule \\
DuckDB & analytics speed & custody of what left & welcome inside the
  transform step \\
Audit tables & cheap provenance theater & rewrite resistance & signed,
  witnessed digest chain instead \\
Public logs & append ordering, equivocation & truth of what was logged &
  bind source/policy/result, then log \\
Hash-the-CSV & transit integrity & provenance, policy, time & the capsule
  is the finished form of this idea \\
\bottomrule
\end{tabular}
\end{center}

% =====================================================================
\section{Advantages}
\label{sec:advantages}

\begin{enumerate}
  \item \textbf{Proof-carrying exports.} Every release carries source
    snapshot roots, lineage closure, policy digest, execution appraisal,
    result root, and a transparency receipt -- one content-addressed object.
  \item \textbf{Offline, counterparty-run verification.} The auditor runs
    the verifier with no network and no call-back to the producer; exit
    code 0 means every named check passed.
  \item \textbf{Fail-closed by construction.} Missing dependencies never
    roll up to PASS; \code{INDETERMINATE} is not a release signal; typed
    errors (\code{E\_*}) replace crashes everywhere, including the WAL and
    page layers.
  \item \textbf{Tamper evidence with names, not vibes.} A mutated leaf,
    swapped purpose, stale policy, replayed grant, or equivocating witness
    each fails a specific named check; the attack gallery asserts all of
    them stay blocked.
  \item \textbf{Rewrite resistance.} The capsule digest covers every sealed
    field; tests mutate each field individually and require the digest to
    move. History rewrite requires forging SHA-384 or forging a checkpoint
    signature -- and witnesses make even a signed fork detectable.
  \item \textbf{Custody across systems.} Lineage edges are bound into the
    capsule and checked acyclic; the same evidence story extends to
    Postgres-sourced snapshots through connectors.
  \item \textbf{Replay-proof releases.} One-use, digest-bound grants: a
    consumed grant cannot be replayed, and a grant for another capsule
    does not apply.
  \item \textbf{Algorithm agility without overclaiming.} Classical and PQ
    transport, signature algorithms -- all labeled per record, so the
    deployment can upgrade primitives without the documentation rotting
    into marketing.
  \item \textbf{Runs where the evidence is discussed.} The identical
    engine -- pages, WAL, MVCC, planner -- runs in Node, in the browser on the
    site's five pages, and on the desktop through the wire protocol, so a buyer can execute the proof claims on the
    site's SQL IDE and the PyQt6 Studio itself rather than trusting a demo video.
  \item \textbf{Auditable surface.} Zero production npm dependencies (the
    single crypto dependency, @noble/post-quantum, is a build-time
    devDependency), no eval/dynamic code, ASCII-only sources, and a
    test suite that runs fully offline. A reviewer can read the whole
    system.
\end{enumerate}

% =====================================================================
\section{Verification and evaluation}
\label{sec:verification}

\subsection{Implementation constraints, enforced by tests}
\label{sec:constraints}
A dedicated suite walks the source tree and fails the build on: \code{eval}
and \code{Function} constructors, dynamic \code{require} or dynamic
\code{import()} in engine code, WebAssembly, non-relative runtime imports in
engine code, any production dependency in \code{package.json}, and
non-ASCII bytes in any source file. These are injection-surface and supply
chain constraints, checked mechanically rather than promised.

\subsection{Test inventory (204 at time of writing)}
Lexer and parser (grammar, spans, error positions), tuple/codec round-trips
and bounds, page encode/decode with CRC failure paths, B$+$tree under random
workloads, memory and encrypted pagers, engine semantics (constraints,
joins, aggregates, planner), MVCC and transaction tests, crash matrix
(adversarial WAL tear offsets), WAL rotation, protocol loop (framing,
auth, idempotency, production-mode refusals, catalog introspection under
role visibility), hybrid PQ handshake (both parties agree; downgrade
attempts diverge), evidence pipeline end-to-end (capsule digest coverage
per field; tamper detection), known-answer vectors, attack gallery
($\geq 10$ attacks, all blocked), the LeetCode-183 demo proof, the Python
client round-trip against the framed server, and the desktop IDE smoke
harness that pattern-matches the Qt syntax highlighter against the
canonical keyword surface. The Windows WebCrypto fallback suite pins the
honest-labeling behavior on runtimes that lack X25519/Ed25519.

\subsection{Browser portability findings}
The same pipeline runs in the browser across the site's five pages (home,
SQL IDE, SECP lab, specs, docs). Two browser-specific
defects were found and fixed by that harness: (i) WebCrypto algorithm
availability differs from Node -- solved by the probe-and-label fallbacks of
Section~\ref{sec:crypto}; (ii) \code{CompressionStream} readable sides have
a zero high-water mark in Chromium, so awaiting \code{write()} before
attaching a consumer deadlocks -- the MDL scorer now attaches the consumer
first. Both classes of bug are invisible to Node-only CI, which is why the
browser run is part of the verification story, not decoration.

A third defect was caught by the \emph{desktop} harness, not the browser:
the wire server's CATALOG frame (type 8) introspected the connection's
pre-AUTH anonymous session instead of opening a session with the authed
identity, so the schema explorer always saw an empty catalog under an
identityless skeleton. The Python round-trip exposed it, the fix routes
introspection through \code{connect(identity)}, and the catalog now reports
the caller's roles and only the tables that identity can see.

% =====================================================================
\section{Limitations (read this section)}
\label{sec:limitations}

\begin{itemize}
  \item \textbf{Single writer.} The engine takes one writer thread; no
    distributed consensus, no sharding. Target scale is bounded regulated
    pipelines, not hyperscale OLTP.
  \item \textbf{Synthetic demo estate.} The site's lab estate (48 opaque
    member rows, three seeded disagreements) is labeled SYNTHETIC
    throughout; capture lag and p95 verifier latency are targets, not
    measurements on a real pipeline.
  \item \textbf{Simulated KMS.} Key custody is envelope keys over real
    AES-256-GCM in process; production HSM/KMS integration is specified,
    not shipped.
  \item \textbf{Simulated witnesses.} Witness acknowledgements are
    HMAC-based stand-ins; quorum checking and split detection are real,
    the federation is not.
  \item \textbf{Host-local time.} Timestamps trust the host clock; a
    trusted timestamping service is required for adversarial settings.
  \item \textbf{Transport-only PQ.} No PQ signatures; ML-KEM protects
    sessions, not artifacts.
  \item \textbf{MDL is a heuristic.} The compression residual is an
    anomaly score, explicitly not Kolmogorov complexity.
  \item \textbf{Verifier trusts its own code.} A verifier checks evidence
    correctly if its binary/source digest is pinned externally; the
    self-digest is printed, the pinning ceremony is the adopter's.
  \item \textbf{Not SQL-compatible.} SEAMQL is deliberately small; porting
    arbitrary SQL is out of scope for 0.1.
\end{itemize}

% =====================================================================
\section{Reproducing everything}
\label{sec:repro}

\begin{lstlisting}[language={},basicstyle=\ttfamily\small]
npm install           # devDependencies only; no production deps
npm run engine:test   # 204 tests: parser..protocol..gallery..drivers..ide
npm run verify        # offline capsule verifier, exit 0 = full PASS
npm run attacks       # adversarial gallery, all must be BLOCKED
npm run vectors       # known-answer vectors
npm run dev           # five-page site: home, IDE, lab, specs, docs

# wire protocol + clients
npx tsx src/server/tcp.ts --mem --port 7301
python - <<'PY'
import sys; sys.path.insert(0, 'drivers/python')
import seamdb
with seamdb.connect('127.0.0.1', 7301, user='root', roles=['owner']) as db:
    db.query("spawn table t: id: whole number, main key, name: text, needed")
    db.query("drop in t { id: 1, name: 'Ada' }")
    print(db.query("pull name from t when id = 1").pop().rows)
PY

# desktop IDE (PyQt6 -- a window, not a web frame)
python ide/seamdb_ide.py --host 127.0.0.1 --port 7301 \
                         --user root --roles owner
\end{lstlisting}

Every command runs offline. The LeetCode-183 script shown in
Section~\ref{sec:seamql} is executed verbatim by
\code{tests/demo-proof.test.ts} and by the browser IDE, so the document, the
page, and the engine cannot drift apart silently.

% =====================================================================
\section{References}
\label{sec:refs}

\begin{itemize}
  \item RFC 8785 -- JSON Canonicalization Scheme (JCS).
  \item RFC 9162 -- Certificate Transparency Version 2.0 (log shape,
    inclusion/consistency proofs).
  \item RFC 8032 -- EdDSA: Ed25519. \quad RFC 7748 -- X25519.
  \item RFC 8446 -- TLS 1.3 (hybrid group design reference).
  \item RFC 5869 -- HKDF. \quad NIST SP 800-38D -- AES-GCM.
  \item FIPS 203 -- ML-KEM (module-lattice-based key-encapsulation).
  \item FIPS 186-5 -- Digital Signature Standard (ECDSA, EdDSA).
  \item RFC 9334 -- RATS architecture (appraisal reference model).
  \item C. Mohan et al., ``ARIES: A Transaction Recovery Method Supporting
    Fine-Granularity Locking \ldots'', ACM TODS 1992.
  \item J. Rissanen, ``Modeling by Shortest Data Description'' (MDL).
  \item @noble/post-quantum -- referenced ML-KEM-768 implementation.
\end{itemize}

\end{document}
